\section{Cohorts, prediction protocol and statistics}
\label{app:setting}

\subsection{Cohorts}
\label{sec:cohorts}

Table~\ref{tab:cohorts} lists the cohorts. They differ in three ways that matter: the sensor interval, the observation
window and the source of the carbohydrate values. We keep subjects with at least ten meals, or five in the cohort with
standardized meals. The Shanghai sensor samples every $15$ minutes and the stored $5$-minute grid was filled by
interpolation, so every trace objective reads only the real samples.

\begin{table}[t]
\centering\small
\begin{tabular}{lrrcccr}
\toprule
cohort & subjects & meals & meals per subject & sensor (min) & window (min) & median carbs (g) \\
\midrule
CGMacros & 45 & 1640 & 36 (17--64) & 5 & 180 & 56 \\
Shanghai & 87 & 2608 & 30 (10--56) & 15 & 180 & 33 \\
Standardized meals & 15 & 82 & 5 (5--6) & 5 & 145 & 35 \\
\bottomrule
\end{tabular}
\caption{The cohorts after the selection rule (at least ten meals per subject, five for the standardized-meal cohort). Meals per subject: median (minimum--maximum). The Shanghai sensor samples every fifteen minutes; its stored five-minute grid is interpolated and trace objectives read the real samples only. Carbohydrate sources are described in the text.}
\label{tab:cohorts}
\end{table}


CGMacros records carbohydrate in a photographed, nutritionist-coded log. The Shanghai cohort has a free-text dietary
record, and we estimate carbohydrate from the food name and weight by keyword. That estimate has its own error, which is
one reason for the carbohydrate-error experiments of Section~\ref{sec:si}. In the standardized-meal cohort the meals are
known by construction. The window is the observation window of the iAUC, and it is shorter in the last cohort.

\subsection{Prediction protocol}
\label{sec:cv}

Held-out prediction uses $5$-fold cross-validation within subject, repeated $5$ times. Fold assignments are seeded by
SHA-256 and identical for every method, so no method benefits from a favorable split. One unit is a subject and a
repeat, every meal is held out once, a subject's score is its mean over repeats, and every interval resamples
\emph{subjects}. We compare gradient fitting with three parameters or with $\SI$ only, exhaustive grid search of the same
loss (a $20\times8\times8$ grid, or a $41$-point grid), and the same under the full-trace objective. The baselines are
the personal mean, the population default, and persistence. Two methods are called equivalent when the $90\%$ bootstrap interval of their paired difference
lies inside $(-150,+150)$ mg\,dL$^{-1}$\,min. Margins of $90$ and $300$ are also reported. A difference is not called an
equivalence unless that test was run.

\paragraph{The margin.}
The equivalence margin of $150$\,mg\,dL$^{-1}$\,min was chosen prior to analysis as an operational scale rather than a clinical cutoff. It is smaller than the gain of the
gradient fit over a subject's own mean on CGMacros (\resDiffGradThreeVsPersonalMean{}, roughly $65\%$ of that gain), and a small fraction of the mean
error (\resMaeGradThree{}). The smaller margin of $90$ is the stricter test and $300$ the looser one, and the
conclusions for the primary comparisons are those reported in Appendix~\ref{app:results}. A margin
tied to a clinical decision would require an outcome beyond iAUC.

\subsection{Statistics and the analysis plan}
\label{sec:plan}

Fractions of subjects carry Wilson intervals. Paired differences carry subject-level BCa bootstrap intervals
($\resBootstrapResamples$ resamples), Wilcoxon tests with Holm correction within families, and two one-sided tests for
equivalence. Ten hypotheses with their thresholds were fixed in version control before the analyses were run
(Table~\ref{tab:plan}); every other analysis is exploratory. The full record, with each later analysis and when it was
declared, is kept with the code.
